% !TeX spellcheck = en_US
% !TeX encoding = UTF-8
\chapter*{Abstract}

In the second lab of the Security Insider Lab, we detect network attacks from flow features with machine learning. We first recorded our own traffic and converted it into flows with CICFlowMeter. Our first captures were mostly IPv6 and produced bogus flows, because CICFlowMeter force-parses non-IPv4 frames as IPv4. With IPv6 disabled, we obtained 569 flows, 547 of them valid TCP/UDP flows. As training data we use the Friday-afternoon DDoS file of CIC-IDS2017, which has 225,745 flows, 128,027 of them DDoS. Pre-processing removed infinite values, duplicate rows and constant and identical columns, which left 223,082 flows with 59 features. We then made a stratified split, min-max scaled the features and selected the 20 best with $\chi^2$. We trained Logistic Regression and Decision Tree classifiers for the binary target BENIGN vs.\ DDoS. The default Decision Tree is best: on the test split of 66,925 flows it makes 7 false positives and 25 false negatives (DDoS F1 0.9996), whereas Logistic Regression raises about 2,000 false alarms (DDoS F1 0.974). Class weighting changes almost nothing, because DDoS is the majority class. Applied flow by flow to our own benign capture, the Decision Tree classified all 547 valid flows as BENIGN, so it raised no false alarm. Since the capture contains no attack, we cannot measure the detection rate there, and the near-perfect benchmark scores may partly rest on shortcut features such as the destination port.
